\section{引言}

随着工业基础设施规模不断扩大以及设备服役年限增加，电力、石化、冶金、交通和新能源等领域对巡检效率、作业安全性和状态感知能力提出了更高要求。输电线路、桥梁、风电设施、厂房、管廊、储罐以及大型工业装备通常具有结构尺度大、设备分布复杂、运行连续性要求高等特点，部分设施还长期处于高空、狭窄、封闭、强反光、粉尘或强电磁等人员不宜长期作业的环境。传统人工巡检能够依赖作业人员的经验完成现场判断，但在大范围、高风险和难抵达区域内，存在劳动强度高、作业效率受限以及巡检结果受人员经验影响等问题。随着无人机平台、机载传感器和自主控制技术的发展，无人机凭借机动灵活、非接触作业、视角可调以及可搭载多种传感器等特点，逐步成为工业设施巡检的重要技术手段\cite{floreano2015science,kumar2012opportunities,shakhatreh2019survey,kanellakis2017survey}。

从应用发展看，无人机最初主要作为人工巡检的数据采集平台，其作业方式依赖操作者规划航线并在飞行过程中完成目标观测。摄影测量和遥感技术的发展，使无人机能够获取高分辨率影像并完成设施三维重建，为基础设施检测和数字化管理提供了新的数据来源\cite{colomina2014unmanned,nex2014uav}。随后，计算机视觉与深度学习逐渐被引入无人机巡检，使系统能够从采集图像中自动识别目标设备、缺陷和异常状态，并推动无人机巡检由“人工飞行、自动采集”向“自动感知、自动识别”演进\cite{nguyen2018automatic,miao2020uav,zhouwq2024review}。在桥梁等基础设施领域，无人机视觉检测已经能够支持裂缝、剥落和腐蚀等表观病害的自动识别\cite{morgenthal2014quality,ham2016visual,dorafshan2018bridge,spencer2019advances}；在受限空间和地下设施中，研究者开始进一步关注自主导航与环境建图问题，例如 Ozaslan 等将自主导航和三维建图用于压力钢管及隧道检测，Petrl{\'\i}k 等则验证了受限环境下无人机自主运行的可行性\cite{ozaslan2017autonomous,petrlik2020robust}。

然而，复杂工业场景下的自主巡检与普通室外环境中的无人机飞行并不等价。工业设施通常具有钢结构密集、几何结构重复、局部空间狭窄、遮挡严重以及光照条件不稳定等特征，GNSS 信号还可能受到建筑物遮挡和多径效应影响；在部分作业区域，粉尘、烟雾、雨雾、强反光和动态人员车辆会进一步降低视觉或激光传感器的观测质量\cite{scaramuzza2014vision,lvyang2019sense,petrlik2020robust}。与此同时，巡检对象本身也并非单一目标：无人机既要识别设备及其部件，还需要发现裂纹、腐蚀、缺失、过热等不同类型异常，因此需要同时获得几何、视觉、运动和状态等多源信息。对于大范围工业设施，单架无人机还受到续航、通信和机载计算资源限制，多无人机协同又引入任务分配、路径冲突、通信和能量管理等问题\cite{zeng2016wireless,mozaffari2019tutorial,zhang2019energy}。

由此可见，复杂工业场景无人机巡检并不是若干独立算法的简单组合，而是一个由实际巡检任务牵引的连续过程。无人机首先需要感知自身状态和周围环境，在此基础上完成可靠定位与地图构建；随后需要针对巡检对象获取视觉、热红外及几何等信息，进一步识别设备缺陷和异常状态；当感知或检测结果发生变化时，还需要调整观察位置和巡检顺序，并在任务规模较大时协调多架无人机共同完成巡检。因此，复杂工业场景自主巡检可以从“环境感知与信息融合—设备缺陷与异常检测—自主任务规划与多机协同”三个相互衔接的层次进行理解。其中，环境感知是无人机安全运行和后续任务执行的基础，缺陷与异常检测是实现设备状态认知的关键，而规划与协同则决定感知和认知结果能否转化为有效的巡检行为。

基于这一认识，本文不按照具体算法逐项罗列相关工作，而是以复杂工业巡检的实际任务为主线，对国内外研究进展进行梳理。首先，从工业场景与巡检任务出发，分析无人机需要获取的信息以及复杂环境对感知系统提出的要求；随后围绕视觉、激光、惯性及其他传感器的信息互补关系，总结环境感知与多模态融合方法的发展；在此基础上进一步讨论设备异常识别以及任务规划和多无人机协同问题。通过比较不同研究在应用场景、任务目标、技术路线和实验条件上的差异，分析现有研究能够解决的问题及其适用边界，并为后续面向复杂工业场景的自主巡检研究提供问题基础。

\section{复杂工业场景无人机环境感知与信息融合}

复杂工业场景无人机巡检的首要问题不是直接选择某一种定位或识别算法，而是明确“无人机需要在什么环境中完成什么任务，并因此需要获取什么信息”。只有将具体工业场景、巡检任务和信息需求联系起来，才能进一步说明为什么需要视觉、激光、惯性以及红外等不同传感器，也才能解释现有自主感知和多模态融合方法的研究动因。因此，本章首先从复杂工业场景及巡检任务出发，随后讨论环境感知与自主定位方法，最后分析多模态信息融合的发展及其仍存在的问题。

\subsection{复杂工业场景及巡检任务分类}

无人机工业巡检的应用对象具有明显的多样性。从已有研究看，相关应用已经覆盖输电线路、桥梁、风电和光伏设施，以及厂房、隧道、压力钢管等多种场景\cite{nguyen2018automatic,miao2020uav,zhouwq2024review,morgenthal2014quality,dorafshan2018bridge,ozaslan2017autonomous}。尽管这些设施在结构形式、工作条件和缺陷类型上存在差异，但从无人机的作业方式看，可以将复杂工业巡检场景概括为大范围开放或半开放设施、复杂内部空间以及高风险或强退化环境三类。不同场景并非相互排斥，实际工业现场往往同时具有多个特征，因此场景分类的目的并不是建立严格的行业边界，而是从无人机所面对的感知和任务约束出发，为后续方法分析提供统一框架。

第一类是大范围开放或半开放工业设施，例如输电线路、桥梁、风电场以及光伏场站等。该类场景通常具有作业范围大、巡检目标分散、设施几何尺度差异显著等特点。无人机需要在较大的空间范围内完成设备搜索、目标定位和区域覆盖，并在接近目标后获取足以支持缺陷判别的高分辨率信息。输电线路巡检研究已经形成较成熟的应用体系，无人机可以依靠航空平台快速获取线路及杆塔部件影像，再利用视觉算法识别绝缘子、导线和其他关键部件以及相应缺陷\cite{nguyen2018automatic,miao2020uav,zhouwq2024review}。桥梁巡检则更加突出结构表面覆盖和病害定位问题，无人机视觉系统已被用于裂缝等病害的自动检测，同时摄影测量与三维重建为后续结构状态评估提供了空间信息\cite{yeum2015vision,cha2017deep,morgenthal2014quality,ham2016visual,spencer2019advances}。此类研究说明，在大范围场景中，“到达目标”并不是巡检任务的终点，如何在保证覆盖效率的同时获得满足检测要求的观测质量，同样是任务设计的重要内容。

第二类是厂房、管廊、储罐、隧道和压力钢管等复杂内部空间。与开放场景相比，这类环境中无人机受到明显的空间约束，GNSS 信号可能减弱甚至完全不可用，建筑物、设备和管路还会产生遮挡和多径效应，从而使依赖卫星定位的传统导航方式难以单独满足要求\cite{nil2012review,scaramuzza2014vision,ozaslan2017autonomous}。在这类环境中，无人机需要依靠机载视觉、激光和惯性传感器建立局部环境模型，并在持续运动过程中估计自身位置和姿态。例如 Ozaslan 等针对压力钢管与隧道环境开展了自主导航和三维建图研究，说明在 GPS/GNSS 受限环境下，环境建图与自主飞行必须结合考虑\cite{ozaslan2017autonomous}。Petrl{\'\i}k 等针对受限环境设计鲁棒无人机系统，也表明复杂空间中的自主运行需要同时考虑感知、避障和系统鲁棒性\cite{petrlik2020robust}。

第三类是具有明显退化因素的高风险工业环境。这类场景不一定具有单一的几何特征，其关键问题在于传感器观测质量可能随环境状态显著变化。例如，粉尘和烟雾会降低视觉清晰度，局部强反光和弱纹理会削弱视觉特征的稳定性，复杂钢结构又可能造成激光遮挡；强电磁环境以及通信条件变化还可能影响无人机的数据传输与协同运行\cite{lvyang2019sense,zeng2016wireless,mozaffari2019tutorial}。因此，在退化工业环境中，自主巡检系统面临的并不是单一传感器性能不足，而是“观测质量随环境变化”的系统性问题。这也意味着，感知系统需要具备检测异常观测、利用不同传感器之间的信息互补以及在部分传感器性能下降时维持运行的能力。

从上述场景进一步看，无人机工业巡检任务可以概括为三个层次。第一层是面向无人机自身安全的环境感知任务，包括障碍物检测、距离估计、动态目标感知以及自身运动状态估计，其直接目标是保证无人机能够在复杂环境中稳定飞行。第二层是面向巡检对象的目标感知任务，包括设备识别、部件定位以及重点区域定位，为后续缺陷检测提供目标和空间范围。第三层是面向场景理解的任务，包括环境地图构建、设备空间关系建模和巡检结果空间关联，从而为后续任务规划和状态管理提供基础\cite{kanellakis2017survey,nex2014uav,tao2019digital}。三类任务之间存在明显的递进关系：无人机首先需要知道“自己在哪里以及周围有什么”，随后才有可能准确回答“要检查什么”，最终才能依据设备和环境状态决定“下一步如何检查”。

因此，复杂工业巡检中的感知需求不宜简单归纳为“需要高精度定位”或“需要识别缺陷”，而应理解为对几何、视觉、运动和设备状态等多类信息的综合需求。可见光图像通常能够提供丰富的纹理和语义信息，适合设备及表观缺陷识别；红外热像能够反映温度分布，适合发现部分电气或热状态异常\cite{jadin2012recent,usamentiaga2014infrared}；激光雷达能够提供直接的三维几何信息，对光照变化相对不敏感，适合定位建图和结构感知\cite{zhang2014loam,wujq2021review}；IMU 能够提供高频运动观测，在无人机快速运动以及视觉或激光短时退化时发挥重要作用；GNSS 在开放环境中能够提供全局位置基准，但在封闭和复杂钢结构区域存在明显局限。不同信息源的功能差异说明，工业巡检天然具有多模态特征，这也构成了下一节环境感知与信息融合研究的现实基础。

总体而言，已有应用研究已经证明无人机能够在多种工业设施中承担巡检任务，但不同场景对应的核心矛盾并不相同：大范围设施更加关注覆盖效率、目标定位和观测质量；复杂内部空间更加关注无 GNSS 环境中的定位、建图和避障；高风险退化环境则更加关注感知信息可靠性与系统鲁棒性。也就是说，工业巡检感知方法的发展并非单纯追求更高的定位精度，而是在不断变化的场景和任务约束下，提高无人机对自身、环境和巡检对象的综合感知能力。基于这一认识，下一步需要从自主定位、环境建图以及障碍物感知等方面考察不同感知技术的适用性与发展脉络。

\subsection{工业巡检环境感知}

环境感知是连接工业巡检任务与自主飞行执行的基础环节，其核心是利用机载传感器持续估计无人机自身状态、周围环境结构以及与障碍物和巡检对象之间的空间关系。早期无人机自主导航研究主要关注飞行器的制导、导航与控制问题，在依赖 GNSS 的室外环境中，可以通过卫星定位和惯性测量实现基本的位置与姿态估计\cite{kendoul2012survey,cai2014survey}。然而，当无人机进入室内、隧道、厂房以及大型设备内部后，GNSS 不再能够稳定提供位置约束，自主定位便需要更多依赖机载视觉、激光和惯性信息\cite{nil2012review,scaramuzza2014vision}。因此，GPS/GNSS 可用性变化成为推动无人机自主感知方法发展的重要实际因素。

在视觉感知方面，相机具有体积小、成本低和信息维度高等优点，能够同时提供纹理、颜色和语义线索，因此长期以来一直是无人机自主感知的重要传感器。视觉 SLAM 将相机观测与位姿估计和地图构建结合起来，是无人机在 GNSS 拒止环境中实现自主导航的典型路线\cite{durrant2006simultaneous,cadena2016past,scaramuzza2014vision}。其中，ORB-SLAM 通过特征点、关键帧和回环检测建立较完整的视觉 SLAM 框架\cite{mur2015orb}，ORB-SLAM2 将其进一步扩展至单目、双目和 RGB-D 相机，ORB-SLAM3 则在视觉和视觉--惯性模式基础上继续支持多地图场景\cite{mur2017orb,campos2021orb}。这些工作推动了视觉 SLAM 从单一视觉里程计向具有回环检测、重定位和多模式运行能力的通用框架发展。

视觉方法的主要优势在于能够同时提供几何与语义信息，但其性能对环境纹理、光照和运动状态较为敏感。在钢结构密集或纹理重复的工业环境中，局部区域可能难以提取稳定的视觉特征；在快速飞行、强反光或低照度条件下，图像质量进一步下降。因此，实际工业巡检中需要考虑视觉算法的适用环境，而不能仅依据公开数据集上的精度判断其工程能力\cite{scaramuzza2014vision,kanellakis2017survey}。视觉--惯性方法利用 IMU 提供的高频运动信息，可以在相机短时间内受到模糊或遮挡影响时维持状态估计。VINS-Mono 是典型的视觉--惯性状态估计方法，通过联合优化视觉和惯性信息实现鲁棒状态估计\cite{qin2018vins}；Loianno 等进一步将单目视觉和 IMU 用于小型四旋翼的估计、控制和规划，说明轻量化视觉--惯性系统能够满足无人机快速飞行需求\cite{loianno2017estimation}。Faessler 等则进一步验证了基于视觉的自主飞行与实时稠密三维建图能力\cite{faessler2016autonomous}。

与视觉相比，激光雷达直接测量空间几何结构，对光照变化不敏感，因此在结构复杂和视觉退化环境中具有明显优势。LOAM 将激光里程计和地图构建结合起来，通过提取并匹配边缘和平面特征，实现了实时三维激光建图\cite{zhang2014loam}。在此基础上，激光与惯性信息逐渐形成紧耦合融合路线。FAST-LIO2 直接利用原始点云与惯性信息进行状态估计，减少了传统激光特征提取环节，并兼顾较高的实时性和估计精度\cite{xu2022fast}；Faster-LIO 则进一步从稀疏体素和并行计算角度降低实时建图开销\cite{bai2022fasterlio}。这些方法表明，激光--惯性融合已经从“分别处理两类传感器”逐步转向“在统一状态估计框架中直接利用多源测量”。

从工程应用角度看，视觉与激光并非彼此替代。视觉的突出优势是纹理和语义表达能力强，激光的突出优势是几何结构稳定，IMU 则提供高频运动约束。对于开放场景和纹理丰富环境，视觉方法具有成本和信息量方面的优势；对于隧道、管廊以及大型钢结构内部，激光及激光--惯性融合方法能够降低光照变化带来的影响；对于快速运动平台，视觉--惯性融合可以借助 IMU 改善短时间运动估计\cite{qin2018vins,xu2022fast,gupta2022slam}。因此，环境感知方法的发展可以概括为“单一传感器感知—惯性辅助—多传感器紧耦合”的逐步演进，而其根本驱动力是单一感知方式在特定工业环境下存在不可避免的性能边界。

与此同时，工业巡检环境感知并不只包括定位和建图，还需要面向飞行安全进行障碍物及动态环境感知。无人机感知与规避研究已经从静态障碍物检测逐步扩展到复杂环境中的在线感知和避障问题\cite{lvyang2019sense}。对于工业厂区而言，设备、管路和结构件通常构成复杂静态障碍物，而人员、车辆及其他移动设备则形成动态障碍物；这要求环境模型不仅能够描述静态几何结构，还要能够反映短时间内发生的环境变化。因此，传统一次性建图与规划方式在动态工业现场存在局限，环境感知系统需要持续更新地图和状态估计结果，并为后续路径规划提供实时信息。

近年来，SLAM 研究进一步向鲁棒感知和主动感知方向发展。Cadena 等指出，SLAM 研究正在从单纯关注定位与建图精度转向更加重视复杂环境中的鲁棒感知问题\cite{cadena2016past}；Gupta 和 Fernando 对无人机 SLAM 与数据融合进行了综述，认为视觉、激光和惯性等多源信息融合是提高 UAV 自主感知可靠性的重要方向\cite{gupta2022slam}。进一步地，主动 SLAM 研究开始关注无人机应该“去哪里才能获得更有价值的信息”，即将运动决策与定位、建图过程联合起来\cite{placed2023active,bircher2016receding}。这与工业巡检的任务特性高度一致：无人机并不是为了建立完整地图而飞行，而是为了完成目标搜索、设备观测和异常确认，因此观测价值本身应成为感知和运动决策的重要依据。

综合来看，现有环境感知研究已经形成较为完整的技术体系：GNSS 提供开放场景下的全局位置约束，视觉 SLAM 和视觉--惯性方法适用于低成本自主导航，激光 SLAM 和激光--惯性方法增强了复杂几何环境下的空间感知能力，而主动 SLAM 开始进一步考虑观测价值与运动决策之间的关系。然而，这些方法在复杂工业场景中仍面临三个共性问题。其一，不同传感器在不同退化条件下的性能差异较大，单一感知模态难以长期保持稳定；其二，多传感器系统对时间同步、空间标定和机载算力具有更高要求，实际部署成本随融合程度增加；其三，现有研究仍有相当部分以“定位精度”或“建图精度”为主要评价目标，与实际巡检任务中的目标观测质量和异常识别需求之间存在脱节。由此，多模态信息融合不应仅被视为提高 SLAM 精度的技术手段，还应承担连接环境理解、设备识别和巡检任务的作用。

\subsection{多模态信息融合}

多模态信息融合的根本原因在于不同传感器对同一环境提供了不同性质的观测。对于无人机工业巡检，可见光相机能够提供颜色、纹理与丰富语义，激光雷达能够提供较稳定的三维几何结构，IMU 提供高频运动变化，GNSS 能够提供全局空间参考，而红外传感器则能够反映温度场信息\cite{jadin2012recent,usamentiaga2014infrared,wujq2021review}。这些信息并非完全重复，而是具有明显互补性。因此，当任务要求同时获得环境几何、无人机状态、设备类别以及设备热状态时，依赖单一传感器往往难以满足全部需求，多模态融合由此成为复杂工业场景自主巡检的重要技术路线。

从信息处理层次看，多传感器融合通常可以分为数据级、特征级和决策级三个层次。数据级融合尽可能保留传感器原始测量信息，但要求较高的时间同步和空间标定精度；特征级融合首先提取不同模态的特征，再在统一表示空间内进行联合估计；决策级融合则在不同传感器分别完成识别或状态估计后，再对结果进行综合判断\cite{panq2003basic,wangj2004survey,hey2000multi}。在无人机自主感知中，早期融合更多采用松耦合方式，即先分别得到视觉里程计、激光里程计或惯性估计结果，再通过滤波或优化进行组合。随着计算能力提高以及紧耦合优化方法的发展，融合过程逐渐向统一状态估计框架转变，使不同传感器能够共同约束无人机位姿和地图状态。

视觉--惯性融合是这一发展过程中的典型路线。VINS-Mono 通过紧耦合优化视觉特征和 IMU 信息，实现了对无人机运动状态的联合估计\cite{qin2018vins}。其研究意义并不局限于提高单次位姿估计精度，更重要的是利用不同传感器的时间尺度互补：相机提供相对丰富的空间信息，但频率和运动模糊会限制其短时响应；IMU 提供高频运动信息，但长期积分会产生漂移。二者结合后能够在各自不足的时间尺度上形成补偿。对于工业巡检这种具有快速运动、局部遮挡和短时视觉退化的任务，这种互补关系具有明确的工程意义。

激光--惯性融合则更加突出几何信息与高频运动信息的互补。FAST-LIO2 直接将原始激光点云与 IMU 信息纳入状态估计，通过增量式地图和紧耦合估计实现较高频率的激光--惯性里程计\cite{xu2022fast}；Faster-LIO 进一步降低了高频建图过程中的数据结构和计算开销\cite{bai2022fasterlio}。这类方法尤其适合工业设施内部的三维几何感知，因为厂房、隧道、管廊和设备结构往往具有丰富的空间边界，即使纹理不足，也能够利用几何信息进行定位。但是，激光系统对设备成本、安装体积和机载计算资源的要求相对较高，因此在轻量化无人机平台上仍需要综合考虑任务收益和系统代价。

多模态融合进一步发展后，融合对象已经不再局限于“定位所需的传感器”，而开始扩展到语义和设备状态信息。一方面，视觉图像提供了设备类别、部件和缺陷的语义信息，激光点云能够提供设备的几何形状和空间位置；两者结合能够为场景中的设备建立更完整的空间语义描述。另一方面，红外热像可以提供普通 RGB 图像无法直接获得的温度信息，可用于辅助识别电气设备异常或其他热状态变化\cite{jadin2012recent,usamentiaga2014infrared}。因此，从工业巡检角度看，多模态融合的价值不仅是让无人机“定位得更准”，还在于让系统能够同时知道“设备是什么、位于哪里以及当前状态如何”。

随着视觉大模型、语义分割和三维场景理解方法的发展，环境感知与语义认知之间的边界也正在减弱。数字孪生研究强调将物理设备、运行状态和数字模型建立关联，为工业设施状态管理提供统一的信息载体\cite{tao2019digital}。在这种思路下，无人机获取的图像、点云和红外数据不应只是一次性巡检记录，而可以进一步被映射到设备或空间对象上，从而形成具有时间和空间关联的巡检信息。这为后续的缺陷检测、历史变化分析以及任务规划提供了基础，也说明多模态融合正在由“传感器融合”逐步向“任务信息融合”发展。

主动感知进一步强化了这种任务导向。传统感知系统通常假设无人机按照既定轨迹获取信息，而主动 SLAM 与下一最佳视点研究开始让无人机根据当前地图不确定性和信息增益选择新的观测位置\cite{placed2023active,bircher2016receding}。对于工业巡检，这一思想尤其重要。例如，当系统仅依据远距离图像发现疑似缺陷时，下一步不一定是继续沿预设航线飞行，而可能需要调整无人机位置和视角，从多个方向获取图像或点云，以降低误检和漏检风险。此时，感知结果本身开始参与运动决策，多模态信息也不再是单纯的输入，而成为后续任务规划的状态变量。

尽管多模态融合已经取得大量研究成果，但在复杂工业场景中仍存在较明显的工程瓶颈。首先，不同传感器的测量频率、噪声特性和坐标系不同，时间同步与空间标定误差会直接影响融合结果；其次，不同模态的退化方式并不一致，例如视觉可能受到光照和烟雾影响，而激光可能受到遮挡和反射特性影响，当某一传感器进入异常状态时，传统固定权重融合方法未必能够及时降低其影响；再次，多模态系统通常需要更高的算力和数据带宽，这对无人机平台的续航能力和边缘计算能力提出更高要求\cite{gupta2022slam,zeng2016wireless,zhang2019energy}。因此，真正适用于工业巡检的多模态融合不仅要追求更高精度，还应具有根据场景变化判断“哪些信息可信、哪些信息应当被削弱以及哪些新信息需要主动获取”的能力。

从现有研究的发展脉络看，多模态感知大致经历了由单一传感器、传感器互补、紧耦合状态估计到任务驱动融合的演进过程。早期研究关注如何通过视觉或激光完成定位与建图，随后利用 IMU 等高频信息改善状态估计，再进一步将视觉、激光与惯性信息在统一框架中融合。当前研究开始关注语义地图、主动感知以及面向具体任务的信息选择，但多数方法的评价仍以定位和建图指标为主，对工业巡检中特有的设备识别、异常检测、观测质量以及巡检任务闭环关注不足。由此可见，复杂工业场景无人机的下一步研究重点并不是简单增加传感器数量，而是根据具体巡检任务构建与环境退化相适应的信息融合机制，并使感知结果能够真正服务后续的设备状态认知和自主巡检决策。

\vspace{0.5em}

综上，复杂工业场景中的无人机环境感知已经从单一 GNSS 或视觉导航逐步发展到视觉--惯性、激光--惯性以及更广义的多模态融合体系。场景和任务的变化是推动技术演进的主要因素：开放环境要求高效覆盖与全局定位，受限空间要求在 GNSS 不可用条件下完成自主定位与建图，而高风险退化环境则进一步要求系统能够识别和应对感知不确定性。与此同时，多模态融合正在从“提高定位精度”的局部目标向“支撑环境理解和任务执行”的系统目标演进。现有研究虽然已经为复杂工业巡检提供了重要技术基础，但感知系统与设备异常认知、任务规划之间仍存在明显的信息断裂，这也为后续从“环境感知”进一步走向“缺陷与异常检测”以及面向任务的闭环自主巡检留下了研究空间。

% 第三章：工业设备缺陷与异常检测
% 说明：本章引用全部来自用户现有 refs(1) (1).bib，未新增未经核对的文献。

\section{工业设备缺陷与异常检测}

复杂工业场景中的环境感知解决的是无人机能否稳定获得自身与周围环境信息的问题，而工业巡检的最终目的并不止于“看见”环境，更在于从获取的图像、点云以及其他传感器数据中识别设备状态和异常。对于电力、交通基础设施、工业厂房以及其他典型巡检对象，无人机已经能够利用机载视觉平台完成输电线路、桥梁和工业设备的表观缺陷识别，相关研究证明了无人机图像在裂纹、腐蚀、部件异常等任务中的应用可行性\cite{nguyen2018automatic,miao2020uav,yeum2015vision,cha2017deep,liu2020image,tao2020detection}。然而，与实验室中的通用目标识别不同，工业巡检中的异常具有类别多、尺度差异大、样本稀缺以及场景变化显著等特点，导致“能够检测”与“能够稳定地完成工程巡检”之间仍存在明显差距。

从任务角度看，工业异常认知至少包含三个层次：首先是判断图像或局部区域是否存在异常，即图像级异常分类；其次是确定异常发生的位置及其类别，即目标级缺陷检测；再次是进一步获得缺陷区域的精细边界，即像素级缺陷分割\cite{taox2021survey}。三类任务分别对应不同的巡检需求：当巡检对象数量多、只需要快速筛查时，图像级判断具有较高效率；当需要定位具体部件或缺陷位置时，需要进行目标检测；当需要评估裂纹、腐蚀和剥落等缺陷的形状、面积或空间范围时，则需要进一步进行分割。随着公开工业异常数据集的发展，工业异常检测开始形成更加标准化的评价体系，其中 MVTec AD 为无监督工业异常检测提供了具有代表性的公开数据基础\cite{bergmann2019mvtec}。因此，本章不按照具体网络名称简单罗列研究，而是围绕“已知缺陷的可靠检测”和“未知或稀缺缺陷的异常认知”两类实际问题梳理方法发展，并讨论这些方法面向复杂工业巡检仍然存在的局限。

\subsection{基于视觉与多模态信息的工业缺陷检测}

工业缺陷检测最早主要建立在人工特征和规则判断基础上。随着深度学习技术的发展，检测方法逐渐由依赖人工设计特征转向通过数据学习目标的层次化表示，检测精度、鲁棒性和适应性由此得到明显提高\cite{taox2021survey}。对于无人机巡检而言，可见光图像具有获取方便、空间分辨率高和语义信息丰富等优点，因此长期以来一直是缺陷检测的主要数据来源。输电线路巡检研究已经形成从人工图像分析向自动视觉检测发展的技术路线，桥梁和建筑基础设施领域也开展了大量基于无人机图像的裂缝和表观病害识别研究\cite{nguyen2018automatic,miao2020uav,yeum2015vision,cha2017deep,liu2020image}。

在监督式目标检测的发展过程中，两阶段方法首先通过候选区域生成与分类回归实现目标定位。Faster R-CNN 在区域提议网络的基础上将候选区域生成过程与检测网络结合，为复杂背景下的目标检测提供了具有代表性的端到端框架\cite{ren2015faster}。对于工业巡检而言，两阶段方法能够较充分地利用区域特征，因此适合对设备部件和缺陷目标进行较精细的定位，但其计算过程相对复杂，在机载平台上的实时运行能力受到算力和能耗约束。随后，以 YOLO 为代表的单阶段检测方法直接在图像上完成目标位置和类别预测，在保持较高检测精度的同时显著提高了推理速度\cite{redmon2016you,redmon2017yolo9000,redmon2018yolov3,bochkovskiy2020yolov4,wang2023yolov7}。SSD 等方法同样采用单阶段思想，通过多尺度特征图完成不同尺寸目标的检测，为实际巡检场景中的快速推理提供了重要基础\cite{liu2016ssd}。

工业缺陷通常具有“小目标、低对比度、形态细长或局部纹理相似”等特点，直接将通用目标检测器应用于巡检图像往往难以获得稳定效果。针对类别不平衡和难样本问题，Focal Loss 通过降低易分类样本对训练过程的贡献，提高模型对难分类样本的关注，为工业缺陷检测提供了应对长尾和类别不平衡问题的有效思路\cite{lin2017focal}。与此同时，多尺度特征建模成为工业缺陷检测的重要方向，其根本原因在于无人机与巡检对象之间的距离会持续变化，使同一缺陷在不同图像中具有明显不同的像素尺度。因而，工业巡检检测器不仅需要“检测到”，还需要在远距离筛查和近距离精细观测之间保持稳定性能。

随着 Transformer 在视觉任务中的发展，工业检测方法进一步从卷积网络主导的局部特征建模转向利用全局关系进行目标表达。DETR 将目标检测描述为直接的集合预测问题，减少了传统检测框架中复杂的候选框后处理过程\cite{carion2020end}；Deformable DETR 进一步通过可变形注意力缓解高分辨率特征计算开销，并增强多尺度目标建模能力\cite{zhu2021deformable}。与此同时，Vision Transformer 和 Swin Transformer 分别从全局自注意力与层次化窗口注意力出发，为视觉表征提供了新的基础架构\cite{dosovitskiy2021image,liu2021swin,vaswani2017attention}。这些方法对工业缺陷检测的意义主要体现在两个方面：一方面，长距离依赖有助于理解缺陷与设备整体结构之间的关系；另一方面，多尺度和层次化表征可以改善不同尺度目标的识别能力。然而，Transformer 类方法通常具有更高的计算和数据需求，因此直接部署到无人机边缘平台仍需要考虑模型规模、功耗与实时性之间的平衡。

针对无人机平台资源受限的问题，轻量化网络成为工程应用中的另一条重要路线。MobileNet 通过深度可分离卷积降低计算量，EfficientNet 则通过网络深度、宽度和输入分辨率的联合缩放提高模型效率\cite{howard2017mobilenets,tan2019efficientnet}。对于巡检无人机而言，轻量化并不意味着单纯追求更高帧率，而是需要综合考虑飞行续航、机载算力、图像分辨率以及缺陷检测精度。特别是在远距离巡检时，为了保证小目标缺陷仍具有足够像素，需要较高图像分辨率，而高分辨率输入又会提高推理开销，因此检测性能与机载资源之间存在天然矛盾。

在典型工业设施应用中，监督式视觉检测已经取得较成熟的研究成果。Cha 等利用深度学习开展混凝土裂缝检测，说明深度网络能够从复杂表面纹理中提取具有判别性的裂缝特征\cite{cha2017deep}；Yeum 和 Dyke 利用视觉信息开展桥梁裂缝自动检测，为无人机辅助基础设施巡检奠定了技术基础\cite{yeum2015vision}；Liu 等将无人机图像与三维重建结合，用于桥梁构件裂缝评估，使缺陷识别进一步与空间几何信息联系起来\cite{liu2020image}；Tao 等面向航拍场景开展绝缘子缺陷检测研究，体现了深度视觉检测方法在电力巡检中的应用价值\cite{tao2020detection}。这些研究共同说明，在设备类型与缺陷类型相对明确、训练样本能够充分获取的情况下，监督学习已经能够较好完成“已知设备--已知缺陷”的检测任务。

但是，真实工业巡检中的问题往往比公开数据集复杂。首先，无人机拍摄视角随飞行轨迹持续变化，同一缺陷会呈现尺度、姿态以及遮挡条件的变化；其次，工业设施中普遍存在金属反光、复杂背景和重复纹理，容易使缺陷区域与背景产生混淆；再次，不同设备型号和运行环境会造成明显的数据分布差异，训练于单一设备或单一场景的数据难以直接迁移到新的对象上\cite{taox2021survey,zhouwq2024review}。此外，裂纹、腐蚀等缺陷通常只占图像中很小比例，当无人机处于远距离巡检状态时，缺陷像素数量进一步减少，这使传统目标检测器更容易出现漏检。因此，监督式检测的发展虽然提高了“已知缺陷”的检测能力，但其有效性仍较强地依赖于训练数据覆盖范围。

除了可见光视觉信息外，多模态数据开始被用于改善特殊工业异常的识别能力。红外热成像能够直接反映目标表面的温度分布，对于电气设备过热、局部热异常以及光伏组件热斑等问题具有可见光图像无法直接替代的优势\cite{jadin2012recent,usamentiaga2014infrared}。当视觉纹理不足而热状态差异明显时，红外信息可以作为独立的异常判别依据；当设备外观缺陷与温度异常同时存在时，可见光与红外信息又能够形成互补。因此，多模态缺陷检测并不是简单增加传感器数量，而是利用不同模态对同一设备状态的不同表征提高诊断可靠性。

从研究演进看，工业缺陷检测已经由“依赖人工特征的视觉识别”逐步发展到“深度学习驱动的目标检测与分割”，再向“Transformer 和轻量化模型结合的实时检测”以及“多模态辅助的状态识别”扩展。其主要推动因素不是网络结构本身，而是工程任务提出了越来越严格的要求：首先要求识别准确，其次要求检测速度足以支持无人机在线巡检，进一步要求面对尺度变化、复杂背景、跨设备和多源信息时保持可靠性。现有监督式方法已经能够较好解决已知缺陷，但其核心假设仍是“训练数据能够覆盖实际缺陷分布”。一旦实际现场出现新的设备型号、新的拍摄环境或者训练集中从未出现的异常类型，仅依赖监督检测器便难以保证识别效果。因此，工业巡检研究正在从“检测已知缺陷”逐渐转向“识别未知和稀缺异常”，这构成了小样本与无监督异常检测发展的直接背景。

\subsection{小样本、无监督异常检测}

与通用目标检测任务相比，工业异常检测最突出的矛盾是“正常样本丰富而异常样本稀缺”。工业设备通常长期处于正常运行状态，真实缺陷的出现频率低，而且不同设备之间的缺陷表现又存在较大差异。对于一些涉及安全的关键设备，真实缺陷数据甚至难以大量采集，这使得建立覆盖所有缺陷类型的监督训练集十分困难。因此，工业异常检测逐渐形成了不同于传统分类与目标检测的研究范式，即利用大量正常样本学习设备的正常分布，再将偏离正常模式的样本判断为异常\cite{taox2021survey,pang2021deep}。

无监督异常检测的早期研究主要围绕生成和重建展开。变分自编码器通过学习数据的潜在概率分布进行重建，为无监督异常检测提供了生成式基础\cite{kingma2014auto}；生成对抗网络则利用生成器和判别器之间的对抗学习建立正常样本的生成模型\cite{goodfellow2014generative}。f-AnoGAN 进一步利用生成模型隐空间进行快速异常判断，说明生成式表示可以在缺少异常标注的情况下用于异常识别\cite{fangan2019}。另一条路线则不直接重建输入，而是学习正常样本在特征空间中的紧致分布。Deep SVDD 通过将正常样本映射到紧凑的特征区域，使偏离该区域的样本被识别为异常\cite{ruff2018deep}。这些早期方法解决了“缺少异常标签”的基本问题，但其检测结果仍容易受到重建误差、特征表示能力和正常数据分布复杂性的影响。

随着深层视觉特征提取能力增强，基于预训练特征和局部特征分布的异常检测方法逐渐成为工业视觉中的重要路线。MVTec AD 数据集覆盖多种工业对象和纹理类别，为比较不同无监督异常检测方法提供了统一基准，也推动了异常检测由单一案例研究向标准化评测发展\cite{bergmann2019mvtec}。在此基础上，研究逐步转向更加稳定的特征空间比较和记忆库方法：通过从正常样本中提取深层特征，并建立其统计分布或记忆表示，再根据测试样本与正常特征的距离判断异常。PatchCore 采用核心集采样和特征记忆库实现局部异常定位，在工业异常检测中具有较好的检测效果与可解释性\cite{roth2022towards}。这一发展说明，无监督异常检测的核心思想已经由“能否生成正常样本”逐渐转向“能否学习稳定的正常特征分布”。

为了进一步减少异常标注依赖，自监督学习开始通过人工构造异常样本建立训练信号。CutPaste 类方法通过对正常图像进行切割、粘贴等操作构造合成异常，使模型能够在没有真实缺陷标签的条件下学习异常判别特征；DRAEM 则通过模拟异常区域进行重建与判别联合训练，提高异常定位能力\cite{zavrtanik2021draem}。自监督路线的优势在于能够利用大量正常样本构造训练任务，但其隐含前提是“人工合成异常与真实工业异常具有一定相似性”。如果真实缺陷在纹理、几何结构或物理状态上与人工扰动差异很大，模型仍可能产生较强的域偏差。因此，自监督方法缓解了标注困难，却并未完全解决真实未知异常的开放性问题。

工业场景还要求异常检测方法兼顾在线运行效率。EfficientAD 面向工业视觉异常检测中的实时性需求，在保持检测能力的同时强调毫秒级推理延迟，为在线质量检测与边缘计算提供了新的技术路线\cite{batzner2024efficientad}。对于无人机巡检而言，实时性具有特殊意义：无人机在飞行过程中获得的数据具有明显的时序性，如果异常检测必须离线上传后处理，则检测结果无法及时影响当前飞行行为；而当异常检测能够在机载或近端边缘设备上快速运行时，检测结果才有可能用于实时调整观测距离、飞行方向和巡检优先级。因此，异常检测速度并不是独立指标，而是能否形成闭环自主巡检的重要基础。

在进一步的发展中，视觉基础模型为少样本和开放世界异常认知提供了新的可能。CLIP 通过大规模图像--文本对比学习获得具有较强语义迁移能力的视觉表征\cite{radford2021learning}；WinCLIP 将这种视觉语言表征引入工业异常检测，探索在少量甚至零异常样本条件下完成异常识别\cite{jeong2023winclip}。SAM 通过可提示分割机制提供了通用的视觉区域分割能力\cite{kirillov2023segment}，Grounding DINO 则进一步支持开放词汇目标检测，使检测目标不再严格局限于预先固定的类别集合\cite{liu2024grounding}；DINOv2 通过大规模无监督视觉学习获得通用的视觉特征，为跨场景视觉理解提供了新的特征基础\cite{oquab2024dinov2}。这些研究表明，工业异常认知正在从“给定异常类别、依赖大量标注”向“利用通用视觉知识发现少见或未知异常”发展。

然而，需要看到的是，基础模型并不能直接等同于工业异常检测问题的完整解决方案。首先，工业异常往往不是简单的语义类别，而可能表现为设备结构轻微变化、局部腐蚀、温度异常或组合状态变化，单纯利用通用视觉语义并不能保证对所有异常具有物理意义上的解释能力；其次，基础模型通常在大规模通用数据上学习，其视觉表征与工业设备图像之间存在明显域差异，模型在公开数据上的零样本能力并不能直接推导出真实工业现场的可靠性；再次，巡检任务需要的不仅是“有没有异常”，还包括异常位置、程度、历史变化以及是否需要进一步复查。因此，面向无人机自主巡检的异常认知必须进一步融合场景几何、设备语义、多模态传感信息以及历史状态信息。

从工程巡检的角度看，红外等非视觉模态能够为无监督异常检测提供重要补充。可见光图像主要描述设备表观结构和纹理，红外热像则描述温度分布，两者对于异常的敏感性具有互补关系\cite{jadin2012recent,usamentiaga2014infrared}。例如，当设备外观没有明显变化而内部热状态发生异常时，仅依赖 RGB 图像可能难以识别，而红外信息可以提供新的异常依据。进一步将视觉、红外以及设备历史信息结合，可以把异常判断从“单图像偏离正常分布”提升到“多源状态是否一致”的判断。这种思路与数字孪生所强调的物理设备、数字模型和运行状态关联具有一致性，即巡检结果可以不再作为孤立图像存在，而是与具体设备和历史状态建立长期关联\cite{tao2019digital}。

因此，小样本和无监督异常检测的发展体现出明显的阶段性。第一阶段主要解决异常样本缺失问题，以生成式模型和特征分布建模为代表；第二阶段进一步通过预训练特征、记忆库和自监督学习提高异常定位的稳定性与效率；第三阶段则开始利用视觉语言基础模型和通用视觉表征探索少样本、零样本及开放世界异常认知\cite{pang2021deep,roth2022towards,batzner2024efficientad,jeong2023winclip}。总体而言，这些方法逐步降低了对大规模缺陷标注的依赖，但真实工业巡检仍存在三个突出问题：一是设备型号、拍摄距离和环境条件变化导致正常样本分布本身发生变化；二是某些异常具有明显的物理和时序属性，单张图像难以完成可靠判断；三是异常检测结果尚未充分参与巡检任务决策。因此，工业异常检测的研究重点正在由单纯提高检测指标，进一步转向跨场景泛化、多模态状态理解以及与巡检决策相结合的闭环异常认知。

综上，工业设备缺陷与异常检测已经形成从监督式视觉检测到无监督异常检测、再到基础模型驱动的开放世界认知的发展脉络。对于已知设备和已知缺陷，深度目标检测与分割方法已经具有较好的应用基础；对于异常样本稀缺和缺陷类别不断变化的实际现场，无监督、自监督及基础模型方法提供了更加灵活的解决思路。然而，复杂工业无人机巡检的最终目标不是获得单张图像上的最优检测结果，而是让异常认知能够与环境感知和后续巡检行为形成连续的信息链。当无人机发现疑似异常后，系统还需要决定是否靠近目标、是否改变观测角度、是否进行重复观测以及如何调整后续巡检顺序。由此，缺陷检测与异常认知将进一步从“感知结果输出”走向“驱动巡检决策”，并与下一章的自主巡检任务规划和多无人机协同形成自然衔接。


\section{无人机自主巡检规划与多无人机协同}

第2章和第3章分别从环境感知与信息融合、设备缺陷与异常检测两个层面总结了复杂工业场景无人机巡检的信息获取与状态认知研究。由此可见，完成一次有效的工业巡检并不只是“按照预先设定航线飞行并采集图像”，而是需要根据环境结构、设备分布、可观测性以及实时检测结果持续决定“去哪里、从什么位置观察、观察多久以及是否需要再次检查”。因此，在获得环境与设备信息之后，如何将感知结果转化为可执行的巡检任务和飞行轨迹，构成无人机自主巡检从“感知”走向“决策与执行”的关键环节。

与一般点到点导航不同，工业巡检通常具有明显的任务属性。一方面，巡检对象往往是具有几何结构约束的设备或设施表面，需要满足覆盖完整性、观测距离、视角、分辨率和安全间隔等要求；另一方面，巡检任务还受到无人机续航能力、飞行安全、禁飞区域、风场、通信以及任务优先级等实际条件约束。因此，现有研究逐渐由单纯的最短路径规划扩展到覆盖路径规划、视点规划、信息增益驱动的主动感知以及动态重规划等问题\cite{Aggarwal2020,Kumar2023}.

进一步地，当巡检区域扩大、目标数量增加或多个无人机具有不同载荷与续航能力时，单机规划难以同时满足任务覆盖与时间约束，多无人机任务分配、路径协同和通信组织成为新的研究重点。多无人机系统的核心不再只是分别为每架无人机寻找一条可行路径，而是需要在任务收益、路径代价、能源消耗、时序约束和通信能力之间进行整体权衡\cite{Skaltsis2023,Chakraa2023}.

\subsection{自主巡检任务规划与动态路径重规划}

工业巡检路径规划的研究基础来源于传统机器人运动规划。对于已知地图和静态目标，规划问题通常可表述为在满足碰撞约束、动力学约束和任务约束的条件下，为无人机寻找从当前位置到目标位置的可行轨迹。采样型规划方法以概率道路图、快速扩展随机树及其优化扩展为代表，能够处理具有较高维度和复杂约束的运动规划问题\cite{Kingston2018,Orthey2024}。基于图搜索的A*及其改进算法具有搜索稳定、易于加入启发函数等特点，而人工势场、群智能和其他启发式方法则更加重视求解速度与全局搜索能力。针对无人机应用，Aggarwal和Kumar[ Aggarwal2020 ]从代表性、协作性和非协作性路径规划等角度对相关方法进行了总结，指出路径长度、最优性、完整性、计算效率、能量效率和避碰能力需要综合考虑。近年来的综合性研究进一步表明，路径规划已经从静态二维最短路问题逐步扩展至三维空间、动态环境和多约束条件下的在线规划问题\cite{Kumar2025}。

对于工业巡检而言，仅以“到达目标”为目标并不足以描述实际任务。工业设施通常包含管道、梁柱、设备外壳、储罐及复杂连接结构，无人机需要沿目标表面保持适当的相对距离和观测方向，并在有限能源下获得足够完整的视觉或三维信息。因此，覆盖路径规划（coverage path planning, CPP）逐渐成为基础任务之一。传统CPP通常先根据作业区域形状进行分解，再利用往返式、螺旋式或栅格式路径完成区域覆盖\cite{Oksanen2009,Galceran2013,Choset2001}. UAV覆盖规划综述显示，完整覆盖、路径长度、任务时间、障碍物和能量消耗之间往往存在明显的多目标权衡；当区域存在复杂边界、禁飞区或动态障碍时，单纯依赖规则扫描路径难以保持较高效率\cite{Cabreira2019,Kumar2023}. 近年来针对建筑和基础设施巡检的研究开始将目标表面几何结构直接纳入路径生成。例如，基于BIM的建筑外立面无人机巡检方法能够从建筑模型生成目标表面和候选观测位置，并进一步规划满足覆盖要求的巡检路径\cite{Huang2023BIM}. 这类方法说明，工业巡检的路径规划需要从“覆盖空间”转向“覆盖可观测目标表面”。

在目标表面已知的情况下，进一步需要解决的是观测位置和观测方向的确定，即视点规划（view planning）。机器人主动视觉研究表明，视点选择不仅影响覆盖完整性，还决定图像分辨率、遮挡程度和三维重建质量\cite{Zeng2020}. 对无人机而言，视点规划通常需要同时考虑相机视场、目标表面法向、成像距离、无人机姿态及可达性。与固定航线相比，基于视点的规划能够减少无效飞行，使轨迹更加贴合设备几何结构。对于未知或部分未知环境，下一最佳视点（next-best-view, NBV）方法进一步根据当前地图中的未知区域、重建误差或信息增益选择下一观察位置，从而将巡检规划由“按照预先设定路线执行”转变为“根据当前认知不断选择最有价值的观察位置”\cite{Bircher2016,Zeng2020,Alsadik2025NBV}. 2025年的NBV综述进一步将信息增益、覆盖率、不确定性降低和观测成本作为主要评价维度，说明主动视点规划正在从简单几何覆盖向面向认知质量的规划发展[NBVReview2025].

当环境和巡检结果随飞行过程变化时，静态规划得到的整条路线可能无法继续适用。例如，现场可能出现临时障碍物、人员或车辆，部分设备区域可能无法进入，也可能因为视觉退化、遮挡或异常检测不确定性而需要重新选择观测位置。为此，滚动时域规划和在线重规划逐渐受到重视。Papaioannou等[ Papaioannou2022 ]将三维巡检规划表述为受动力学和感知约束的优化控制问题，并采用滚动时域思想生成可执行轨迹；Ma等\cite{Ma2022}指出，传统滚动时域控制通过局部窗口优化提高实时性，但搜索范围有限，面对不确定障碍和多重飞行约束时仍可能陷入局部可行解。面向密集障碍环境的研究进一步将滚动时域、混合整数规划与局部避碰结合，使路径能够在受限感知范围内持续更新\cite{Hu2024}. 这类方法的共同特点是保留较完整的在线约束表达能力，但相应计算负担较大。

近年来，学习型规划方法开始用于处理难以显式建模的环境不确定性。深度强化学习通过状态、动作与奖励机制学习从环境观测到规划策略的映射，能够在复杂环境中快速给出局部决策\cite{Kumar2025}. 在基础设施巡检中，Pan等\cite{Pan2024}进一步考虑不确定信息和人工巡检资源，利用双重深度强化学习生成无人机与人工协同的巡检路线，说明学习型方法已经从纯导航向任务级规划扩展。然而，学习型规划仍受到训练数据分布、仿真到现实迁移、约束可解释性和安全保证等因素影响。与经典优化方法相比，其优势主要表现为在线推理速度和复杂策略表达能力，而工程部署仍需要对安全约束、可验证性和异常状态处理进行额外设计。

综合来看，单机自主巡检规划大体经历了“最短路径规划—区域覆盖规划—目标表面视点规划—感知驱动的主动规划—动态在线重规划”的发展过程。其研究重点逐渐由单一几何意义上的路径最短转向“在有限飞行代价下获得最大有效巡检信息”。但是，目前许多工作仍然默认任务目标和重点区域在任务开始前已经确定，规划器与缺陷识别模块之间的反馈关系较弱。当第3章所述检测系统发现异常目标以后，现有路径规划是否能够主动改变观测距离、视角和重复巡检次数，仍缺乏统一的闭环建模。因此，未来自主巡检规划需要进一步考虑检测置信度、异常严重程度和信息不确定性，使“发现异常”能够真正转化为“主动复查”。

\subsection{多无人机任务分配与协同决策}

当工业巡检范围扩大或任务目标分散时，单架无人机受制于续航、载荷和飞行时间，采用多无人机并行执行巡检任务可以提高覆盖效率和任务冗余度。多无人机任务规划通常包含任务分配、任务排序、轨迹规划和协同避碰等相互关联的子问题\cite{Skaltsis2023,Chakraa2023}. 其中，任务分配决定“由哪一架无人机执行哪个任务”，路径规划决定“以什么顺序和轨迹完成任务”，而协同决策则需要保证不同无人机的选择不会产生冲突。对于工业巡检而言，任务还具有异质性：不同设备可能需要不同传感器，不同缺陷复查任务可能对应不同距离与视角要求，因此无人机并非完全同质，任务分配不能只依据空间距离。

早期多无人机任务分配研究主要采用集中式优化方法，将所有无人机状态和任务信息汇总至中心节点，通过整数规划、混合整数规划、车辆路径问题、旅行商问题或启发式算法寻求整体较优解。这类方法能够显式处理任务优先级、时间窗、能源预算等约束，也容易获得较好的全局性能。对于规模较小且通信可靠的工业巡检任务，集中式方法具有建模清晰、结果可解释等优势。但随着无人机数量和任务数量增加，其组合复杂度迅速增长，重新求解的计算代价可能成为实时应用的瓶颈\cite{Chakraa2023}. 因此，针对大规模任务和动态任务，分布式方法逐渐成为研究的重要方向。

市场机制和协商机制是分布式任务分配中的代表性方法。CBBA等一致性拍卖算法通过“局部竞价—全局一致”机制，在无需中心控制器的情况下形成冲突较少的任务分配结果。相关研究表明，这类方法兼顾了一定的实时性、分布式特征和通信可实现性，适合网络结构变化较快的多智能体系统\cite{Choi2009,Kim2020}. 后续研究开始进一步考虑任务耦合和异构资源，使无人机能够依据自身能力、载荷和任务收益参与竞价。对于大规模系统，聚类分层和分组协商方法能够降低全局通信量，将一个大规模任务分配问题拆解成若干局部问题\cite{Fu2019,Dong2025}. 这一发展说明，多无人机协同正在从“寻找全局最优分配”转向“在有限通信与计算资源下获得可接受且可快速更新的分配”。

工业巡检中的另一个特点是任务不是完全静态的。第3章中的异常检测可能在飞行过程中产生新的重点任务，例如发现疑似裂纹后需要追加近距离观测，发现热点后需要增加红外复查，或者由于设备遮挡导致原计划任务无法完成。因此，任务集合本身会随感知结果变化，形成动态任务分配问题。动态多无人机任务分配研究通常需要在任务到达、无人机故障、任务优先级变化和资源状态变化时进行增量式重分配，而不是每次从零开始求解\cite{Skaltsis2023,Mahto2023}. 2024年的多无人机动态任务分配综述进一步指出，动态事件、通信约束、系统鲁棒性和规模扩展是实际部署中最关键的挑战之一\cite{Alqefari2025}.

在任务执行阶段，任务分配和路径规划之间并非完全独立。传统方法往往先完成任务分配，再单独优化每架无人机的轨迹，但如果真实飞行距离、障碍约束和访问顺序在分配阶段被过度简化，就可能产生理论上合理而实际执行代价较高的任务分工。针对这一问题，一些研究开始联合考虑任务分配与轨迹设计。例如，针对输电线路巡检的多无人机研究同时考虑山地风场、任务分配与飞行路径，通过联合优化降低整体任务成本\cite{Li2024Wind}. 针对大规模电力巡检的研究则采用聚类与轨迹优化相结合的方式，以减少总巡检时间\cite{Gao2024Power}. 这些研究表明，多无人机协同的研究趋势已经从“先分配、后规划”的串行结构逐步向任务与轨迹联合优化发展。

除任务分配外，多无人机协同还涉及通信与信息共享问题。工业现场可能存在遮挡、金属结构、电磁干扰以及通信盲区，所有无人机持续向中心节点发送完整状态和图像信息会产生明显通信负担。因此，分布式协同需要考虑哪些信息必须共享、何时共享以及共享到什么程度。Kim等\cite{Kim2020}通过分组一致性降低CBBA中的通信开销；Dong等\cite{Dong2025}进一步通过聚类将异构无人机划分为局部群组，在群组内部和群组之间采用不同的一致性机制，从而降低大规模任务分配中的通信量。另一方面，任务状态和无人机状态存在延迟或不一致时，协同决策还可能建立在过期信息之上，这会降低任务分配稳定性。因此，面向实际工业环境的多无人机系统需要将通信可靠性、不确定性和任务价值联合纳入协同决策。

强化学习与多智能体强化学习为多无人机协同提供了另一条技术路线。与基于规则或显式优化模型的方法相比，多智能体强化学习可以通过训练学习无人机之间的协同行为，在复杂和动态环境下减少在线组合搜索。然而，多智能体强化学习容易面临联合动作空间快速增长、训练不稳定、奖励设计困难以及现实环境迁移问题。相比之下，优化、拍卖和一致性方法具有较好的可解释性和约束处理能力，但在高度动态条件下往往需要反复执行迭代。由此形成了当前较为明显的研究趋势：利用优化方法保证基本任务可行性和约束满足，再利用学习方法提高复杂场景下的适应能力，或者采用分层结构将全局任务分配与局部自主决策分开处理\cite{Chakraa2023,Pan2024,Dong2025}.

总体来看，多无人机协同研究形成了“集中式优化—分布式协商—异构资源匹配—动态任务重分配—任务与轨迹联合优化”的发展脉络。对于工业巡检而言，现有研究已经能够较好地处理静态目标集合、固定任务收益和相对确定的环境约束，但仍存在一个重要缺口，即多无人机的任务分配通常以“任务已经定义”为前提，而任务本身实际上可能来自实时感知与异常识别结果。也就是说，真正的工业自主巡检不是简单地将预定义任务均匀分给多个无人机，而应该让“感知结果改变任务集合，让任务价值改变资源配置，让执行结果继续影响后续感知”。这一问题要求任务分配、路径规划与异常认知形成统一的闭环决策框架。

\section{总结与展望}

前述研究表明，复杂工业场景下无人机自主巡检已经形成较为完整的技术体系：在环境层面，通过视觉、LiDAR、IMU、GNSS及红外等多源信息获得无人机位姿和场景结构；在设备层面，通过目标检测、分割、异常检测与多模态识别判断设备状态；在决策层面，通过覆盖路径规划、视点规划、在线重规划以及多无人机任务分配实现巡检任务执行。现有研究不断提高单个模块的精度与效率，但从完整工程过程看，“感知—认知—规划—执行”的环节仍存在较明显的割裂，尤其是感知不确定性没有充分传递给后续规划和协同决策。



\subsection{现有研究存在的主要问题}

（1）复杂退化工业环境下的感知可靠性仍不足。第2章表明，不同传感器在复杂工业现场具有不同的失效模式。视觉容易受到光照、反光、遮挡和纹理不足影响，LiDAR可能受到点云稀疏、表面反射与遮挡影响，GNSS在室内、地下或大型设备附近可能无法稳定工作。虽然多模态融合能够利用不同传感器的互补性，但大量方法默认传感器状态正常，较少根据实际信息质量动态改变融合权重。因此，下一阶段需要从“多模态融合”进一步发展到“面向不确定性的自适应信息融合”。

（2）环境感知与缺陷识别之间仍相对独立。现有定位建图模块主要关注无人机自身状态与环境几何，而缺陷检测模块主要关注设备图像中的异常目标。两类信息通常在系统层面进行串联，而不是相互反馈。例如，检测结果的不确定性可以作为再次观测的依据，但多数规划器并不知道检测模型当前是否确信；反之，规划器产生的观测距离和视角又会直接影响缺陷检测性能。因此，有必要建立面向任务的感知机制，使定位、建图、目标识别和异常检测能够共同决定下一次观测动作。

（3）未知异常与跨场景泛化能力仍是工业巡检的重要瓶颈。工业设备的缺陷类别具有明显的长尾特征，新设备、新材料和新工况可能产生训练集没有覆盖的异常。监督检测模型在已有类别上的性能可能很高，但当设备更换、成像条件发生变化或出现新的缺陷类型时，其泛化能力会下降。小样本与无监督方法缓解了标注依赖，却又面临正常样本污染、阈值设定以及异常语义解释等问题。未来需要把异常检测的不确定性和开放集能力纳入自主巡检闭环，而不是仅输出一个固定类别标签。

（4）现有巡检规划仍以“任务预先确定”为主要前提。覆盖规划和任务分配往往假设目标集合在任务开始前已经给定，异常检测结果只是任务完成后的输出。这种结构难以描述真实巡检中的“发现—复查”过程。一旦无人机发现一个疑似高风险异常，最合理的行为可能不是继续执行原航线，而是立即改变视点、缩短观测距离、调用其他传感器或请求另一架无人机复核。因此，巡检规划需要从固定任务路径转向由感知结果持续驱动的任务生成与重规划。

（5）多无人机系统的协同决策仍缺乏感知闭环和工程约束的统一考虑。当前多无人机任务分配已经能够处理任务数量、资源异构、通信开销和动态事件等问题，但多数工作中的任务价值仍是预定义的。对于工业巡检，更合理的任务价值应当动态变化：设备风险越高、检测置信度越低、信息缺口越大，其复查任务优先级就越高。同时，通信延迟、带宽限制、无人机剩余能量、载荷能力及安全距离都应进入统一决策过程。由此，多无人机自主巡检本质上是一个具有不确定感知、动态任务生成和资源约束的在线决策问题。

\subsection{下一步研究方向}

针对上述问题，后续研究可以围绕“面向复杂工业场景的感知—认知—决策一体化自主巡检”展开，并重点关注以下几个方向。

（1）构建面向信息质量的多模态环境感知方法。现有多模态融合通常按照固定结构进行，而真实工业环境中的不同传感器可靠性会随位置、视角、天气和设备表面特性动态变化。因此，可进一步建立显式的信息质量评价机制，根据观测不确定度、遮挡程度、点云完整性和图像质量自适应调整不同模态的贡献，使感知系统具备“知道自己何时不可靠”的能力。在此基础上，可以将定位、三维重建与目标识别联合建模，使地图不仅包含几何信息，还能够保存设备类别、部件状态以及异常区域等语义信息。

（2）发展面向异常发现的主动感知与闭环巡检机制。第3章已经表明，缺陷检测的结果并非最终目的。未来可以把异常检测置信度、正常性评分以及模型不确定性直接引入视点选择和路径规划，使无人机能够主动选择最有价值的观察位置。当一个异常具有较高风险或较低识别置信度时，无人机可以重新调整距离和视角，或者调用不同传感器进行复查。这样可以形成“检测—评估—复查—更新”的闭环过程，把主动视觉和工业异常检测真正结合起来。

（3）建立任务分配、轨迹规划和感知价值联合优化的多无人机框架。传统多无人机任务分配通常以任务完成时间、距离和能耗为主要目标，而未来可进一步将“信息价值”作为任务收益的一部分。例如，对于不同设备和异常区域，可以根据风险等级、检测置信度、历史缺陷频率和当前信息缺口定义动态任务价值。多无人机系统随后根据各自位置、续航能力、载荷类型和通信状态进行任务竞争与协商。这样可以使“哪架无人机去哪里”不再只由距离决定，而由“完成该任务能够获得多少有效巡检信息”共同决定。

（4）研究面向动态与不确定环境的安全在线重规划。工业现场存在人员移动、车辆穿行、临时障碍物和风场变化等因素，要求规划方法不仅能找到一条可行路径，还要在环境发生变化时保持任务连续性。后续可以将滚动时域优化、局部轨迹生成和学习型策略结合，在全局任务层面保持目标一致，在局部执行层面对突发事件快速响应。同时，需要将无人机动力学、最小安全距离、剩余能量和返航约束显式纳入规划，以提升算法从仿真环境向工业现场迁移的可靠性。

（5）加强真实工业场景验证与统一评价。当前大量研究仍以仿真、实验室数据或公开视觉数据集为主要验证手段，而工业现场具有明显的环境差异和长期运行特征。后续研究应进一步建立包含不同设备类型、不同光照和遮挡条件、不同缺陷等级以及动态环境的测试体系，并从单项算法指标扩展到系统级评价。除了定位误差和检测精度外，还应综合评价单位能耗获得的信息量、异常复查成功率、任务完成时间、漏检率、误报率、通信开销和安全风险。通过统一的任务级评价指标，才能真正判断一种自主巡检方法是否具有工程应用价值。

总体而言，复杂工业场景无人机巡检未来的发展重点并不是继续单独提高某一种检测器、定位器或路径规划器的性能，而是解决不同模块之间“信息如何流动、决策如何反馈”的问题。理想的自主巡检系统应形成闭环：无人机首先利用多模态信息建立对环境和设备的认知，根据检测结果识别重点区域和不确定区域，再结合任务价值、剩余能源与环境风险进行路径重规划和多机任务分配；执行结果又反过来更新地图、设备状态和异常判断。由此，研究对象将由单一算法进一步转向“感知—认知—决策—执行”的自主巡检系统。该发展方向既能够承接当前环境感知、异常检测和多无人机规划等研究基础，也更符合复杂工业场景对安全性、实时性、适应性和长期自主运行能力的实际需求。




